\scanpage{5}

% archive sheet number: 4

\[
\left\{\begin{aligned}
&q_U(y).1_C=\unsure{J|I|\eta}_U(y).1_{\deleted{\unsure{C}}C'}\\
&x,y\in C,\ u\in U,\ \Delta_U u=\sum v_i\otimes w_i\quad\Longrightarrow\quad
\unsure{p|\rho}_U(u)(xy)=\sum(p_U(\deleted{\unsure{u}}v_i)x)(p_U(w_i)y)
\end{aligned}\right.
\]

Si enfin $C=C'$, alors les \deleted{\unsure{en ce}} axiomes d'associativité et d'unité signifient que \deleted{$U$ est un} $U\to\operatorname{Hom}_A(C,C)$ est un \uline{hom de $A$-algèbres unitaires}. Considérant \uline{alors $C$ comme un $U$-module} \inserted{\uline{unitaire}} \uline{à gauche}, les deux relations précédentes s'écrivent aussi

\[
\left\{\begin{aligned}
&u.1_C=J(u).1_C\\
&x,y\in C,\ u\in U,\ \Delta_U u=\sum v_i\otimes w_i\ \Rightarrow\
u(xy)=\sum(\unsure{u_i|v_i}x)(w_i y)
\end{aligned}\right.
\]

Nous voulons maintenant exprimer que $C$, considéré comme ayant son groupe d'opérateurs $B$ dans la catégorie $\unsure{\mathcal{C}|\mathscr{C}}$ des $A$-algèbres comm. avec unité, est \deleted{\uline{principal} de base \unsure{homogène}} \inserted{fibré} \uline{principal} \deleted{de base $D$} \inserted{\uline{homogène}} \deleted{($D\unsure{\to|=|\text{est}}$ une $A$-algèbre commutative). Cela signifie donc qu'il existe un \uline{hom de $A$-algèbres unitaires}}

\diagram{\begin{tikzcd}[ampersand replacement=\&]
\deleted{C\otimes_A C}\& \deleted{B}\arrow[l,draw=none,"{\deleted{\longleftarrow}}" description]
\end{tikzcd}}
